% titel: Blatt 8, A4: Phasendiagramm auswerten (Punkte A-G)
% Teil II: der durchdachte Weg.

\subsection*{Gegeben und gesucht}
Ein allgemeines Phasendiagramm (p gegen T) mit markierten Punkten.
Gefragt nach Phase an A und G, Natur der Punkte B und D, Prozessen an E und F,
sowie dem Verhalten des Schmelzpunkts bei Druckerhöhung.

\subsection*{Ansatz}
Die y-Achse ist der Druck, die x-Achse die Temperatur. Die drei Regionen
(fest, flüssig, gasförmig) sind durch die Koexistenzlinien getrennt.

\subsection*{Auswertung}
\begin{itemize}
  \item \textbf{A} liegt in der Region \emph{fest} (linker, oberer Bereich).
  \item \textbf{G} liegt in der Region \emph{gasförmig} (rechts unten bei
        niedrigem Druck).
  \item \textbf{B} ist der \emph{Tripelpunkt} (drei Phasen im Gleichgewicht).
  \item \textbf{D} ist der \emph{kritische Punkt} (Ende der Dampfdruckkurve).
  \item \textbf{E} liegt auf der fest/flüssig-Linie $\rightarrow$
        \emph{Schmelzen} (f$\to$l) bzw. beim Abkühlen \emph{Gefrieren}
        (l$\to$f); beide Richtungen sind möglich.
  \item \textbf{F} liegt auf der flüssig/gasförmig-Linie $\rightarrow$
        \emph{Sieden/Kondensieren} (l$\to$g ggf. umgekehrt).
\end{itemize}

\subsection*{Schmelzpunkt bei Druckerhöhung}
Die fest/flüssig-Linie (durch den Schmelzbereich) verlängert sich bei
steigendem Druck zu \textbf{niedrigeren Temperaturen}: je höher der Druck,
desto niedriger der Schmelzpunkt. (Das ist das typische Verhalten von
Wasser, dessen Eis unter Druck schmilzt.)

\begin{ergebnis}
A: fest; G: gasförmig; B = Tripelpunkt; D = kritischer Punkt;
E = Schmelzen/Gefrieren; F = Sieden/Kondensieren. Bei Druckerhöhung sinkt
der Schmelzpunkt.
\end{ergebnis}

\begin{kontrolle}
Der kritische Punkt ist der rechte Endpunkt der Dampfdruckkurve, der
Tripelpunkt der Schnitt aller drei Linien --- genau so sind die Buchstaben
im Standarddiagramm üblich. Die nach links gekippte fest/flüssig-Linie
bestätigt das Sinken des Schmelzpunkts.
\end{kontrolle}

\begin{merksatz}
Phasendiagramme liest man so: Region = stabile Phase; Linie =
Koexistenz/Umwandlung; Tripelpunkt = Gleichgewicht aller drei Phasen;
kritischer Punkt = Ende der Dampfdruckkurve.
\end{merksatz}
